\section{Modelado matemático}
El problema abordado puede formularse como un modelo de programación lineal entera. El objetivo de este problema consiste en asignar un vértice a un clique, de forma que se cumplan las condiciones de compatibilidad y se minimice el número total de cliques utilizadas.

\subsection{Conjuntos y parámetros}
En esta sección se define los distintos elementos a considerar como básicos en el modelo, el grafo entrada y los indices a gestionar , los cuales se utilizan para representar las posibles cliques.

Sea $G = (V, E)$ un grafo no dirigido, donde:
\begin{itemize}
    \item $V$ es el conjunto de vértices, con $|V| = n$
    \item $E$ es el conjunto de aristas
    \item $K = \{1,2,\dots,n\}$ es el conjunto de posibles cliques (cota superior)
\end{itemize}

\subsection{Variables de Decisión}

Las variables de decisión permiten representar la asignación de los vértices a las cliques y determinar cuáles de ellos se encuentran efectivamente en la solución.

\begin{itemize}
    \item $x_{ik} \in \{0,1\}$: vale 1 si el vértice $i \in V$ es asignado a la clique $k \in K$, y 0 en caso contrario.
    \item $y_k \in \{0,1\}$: vale 1 si la clique $k \in K$ es utilizada, y 0 en caso contrario.
\end{itemize}

\subsection{Restricciones}

Las siguientes restricciones garantizan que la solución represente una partición valida en cliques, la cual respete la asignación de los vértices que dan a la solución y la forma del grafo.

\textbf{1. Asignación única de vértices}

Cada vértice debe pertenecer exactamente a una clique:
\[
\sum_{k \in K} x_{ik} = 1 \quad \forall i \in V
\]

\textbf{2. Activación de cliques}

Un vértice solo puede ser asignado a una clique si esta está activa:
\[
x_{ik} \leq y_k \quad \forall i \in V, \; \forall k \in K
\]

\textbf{3. Condición de clique}

Dos vértices que no están conectados no pueden pertenecer a la misma clique:
\[
x_{ik} + x_{jk} \leq 1 \quad \forall (i,j) \notin E,\; i \neq j,\; \forall k \in K
\]

\subsection{Función objetivo}

Con las variables de decisión se tienen las condiciones de presentar la asignación de los vértices a las cliques y determinar si estas efectivamente aparecen en la solución. 

\vspace{0.5cm}

Minimizar la cantidad de cliques necesarias para cubrir los vértices de un grafo:

\[
\min \sum_{k \in K} y_k
\]

\subsection{Versión de decisión}

La forma de decisión del problema ofrece la posibilidad de contrastar la existencia de una solución restringida a un número máximo de cliques 

Dado un entero $p$, se desea determinar si existe una solución factible tal que:
\[
\sum_{k \in K} y_k \leq p
\]

cumpliendo todas las restricciones anteriores.